\chapter{Decision Metatheory}
\label{app:decision-metatheory}

When is the question whether every admitted fine Gaussian has an admitted
coarse presentation of its exact marginal decidable, and what does a decision
procedure deliver? For families given by finite semialgebraic descriptions,
real quantifier elimination decides it and constructs the exact good and
failure loci (\Cref{thm:cg-semialgebraic-kron-classification}); the same holds
for singular Gaussian laws on a declared quotient reference with variable rank
(\Cref{thm:cg-variable-support-classification}). The decision does not extend
to families given by programs that enumerate them
(\Cref{prop:cg-kron-enumerator-obstruction}); selectors are measurable but in
general neither continuous nor equivariant
(\Cref{prop:cg-kron-selector-boundary}); and a change of rank destroys
statistical regularity (\Cref{prop:cg-rank-change-statistical-obstruction}).
The appendix supplies the last two rows of \Cref{tab:cg-kron-regimes} and the
only treatment of singular Gaussian laws in this supplement; no result of the
main manuscript depends on it. Notation follows
\Cref{ch:gaussian-coarsegraining}; in addition \(x\) and \(y\) denote tuples of
real variables of first-order formulas, \(x\) for the fine family and \(y\) for
a retained presentation.

This appendix claims no complexity bound useful at any practical scale, no
numerical sign tolerance, no decision procedure for parameters known only
through approximations, no closed-form catalog of closed families, and no limit
or fixed-point theorem for iterated compositions. The classification of
families outside finite encodings is open
(\Cref{open:cg-kron-classification-beyond-finite}). \status{NOT-CLAIMED}

\section{Finite semialgebraic classification}
\label{sec:cg-semialgebraic-kron}

\definitionheading{Finite semialgebraic description}{def:cg-semialgebraic-description}
A finite semialgebraic description of a Kron problem fixes finitely many
labeled nodes with finite coordinate spaces and coordinate Lebesgue
references, a finite list of oriented edge copies and loops, and a finite list
of retained coordinate sets. The fine parameters \(x\in\R^p\) range over the set
defined by a finite formula \(\mathrm{Adm}(x)\), a Boolean combination of
polynomial equalities and strict or weak inequalities, possibly with
existentially quantified auxiliary variables, whose coefficients are real
algebraic numbers, each given by a defining integer polynomial and a rational
isolating interval. The fine precision \(\Lambda(x)\) and source
\(\mathsf h(x)\) have entries that are rational functions of \(x\) on an
explicitly specified nonzero-denominator domain, and every admitted
\(\Lambda(x)\) is SPD. A retained presentation is a tuple \(y\in\R^q\) with a
semialgebraic admissibility formula \(\mathrm{Adm}_{\mathrm c}(y)\), rational
coarse precision and source \(\Lambda_{\mathrm c}(y)\) and
\(\mathsf h_{\mathrm c}(y)\), and a finite semialgebraic compatibility formula
\(\mathrm{Comp}(x,y)\). Correlated weights, nonconvex domains, variable anchors,
nonflat links, polynomial path identities, and finite unions of graph
inventories are allowed; singular or rank-changing supports are treated in
\Cref{sec:cg-variable-support-quotients}.

\propositionheading{Real quantifier elimination}{prop:cg-real-quantifier-elimination}
There is an algorithm that transforms every first-order formula in the language
of ordered fields with integer coefficients into a quantifier-free formula in
the same free variables that is equivalent to it in every real closed field. In
particular the truth of every such closed sentence over \(\R\), and over the
field of real algebraic numbers, is decidable, and the two truth values agree.
\status{STANDARD}

This is Tarski's decision method for elementary algebra \citep{Tarski1951};
\citet[pp.~533--562]{BasuPollackRoy2006} treat quantifier elimination
algorithmically, and \citet[Sec.~6]{CohenMahboubi2012LMCS} give a formally
verified proof over real closed fields with decidable comparison. The decision
theorem below uses only termination.

\lemmaheading{Polynomial encoding of the exact Gaussian marginal}{lem:cg-semialgebraic-encoding}
Under \Cref{def:cg-semialgebraic-description}, with \(M=\Lambda_{\mathrm{hh}}(x)\)
for a nonempty hidden coordinate set, put
\begin{equation}
\begin{aligned}
\widetilde\Lambda_{\mathrm r}(x)&=(\det M)\Lambda_{\mathrm{rr}}-\Lambda_{\mathrm{rh}}(\operatorname{adj}M)\Lambda_{\mathrm{hr}},\\
\widetilde{\mathsf h}_{\mathrm r}(x)&=(\det M)\mathsf h_{\mathrm r}-\Lambda_{\mathrm{rh}}(\operatorname{adj}M)\mathsf h_{\mathrm h},\\
(\det M)\Lambda_{\mathrm c}(y)&=\widetilde\Lambda_{\mathrm r}(x),\qquad
(\det M)\mathsf h_{\mathrm c}(y)=\widetilde{\mathsf h}_{\mathrm r}(x),
\end{aligned}
\label{eq:cg-semialgebraic-schur}
\end{equation}
and for an empty hidden set use \(\det M=1\), \(\widetilde\Lambda_{\mathrm r}=\Lambda\),
and \(\widetilde{\mathsf h}_{\mathrm r}=\mathsf h\). Let \(\mathrm{Rep}(x,y)\) be the
conjunction of \(\mathrm{Adm}_{\mathrm c}(y)\), \(\mathrm{Comp}(x,y)\), the SPD
condition on \(\Lambda_{\mathrm c}(y)\), all denominator admissions, and the last
line of \eqref{eq:cg-semialgebraic-schur}. After finitely many auxiliary
variables are introduced, \(\mathrm{Rep}\) and \(\mathrm{Adm}\) are finite
first-order formulas with integer coefficients, equivalent to the originals at
every real parameter. For admitted \(x\), \(\mathrm{Rep}(x,y)\) holds exactly when
\((\Lambda_{\mathrm c}(y),\mathsf h_{\mathrm c}(y))\) are the natural parameters of the
exact normalized marginal of
\(\mathcal N(\Lambda(x)^{-1}\mathsf h(x),\Lambda(x)^{-1})\) on the retained
coordinates and the presentation is admitted; that marginal precision is SPD,
\begin{equation}
z^\top\Lambda^{\mathrm{eff}}z
=\min_u (z,u)^\top\Lambda(z,u)>0,
\qquad
u_{\min}=-M^{-1}\Lambda_{\mathrm{hr}}z,
\qquad z\neq0 .
\label{eq:cg-semialgebraic-schur-positive}
\end{equation}
\status{ESTABLISHED}

\noindent\emph{Proof.}
The adjugate identity \(M^{-1}=\operatorname{adj}M/\det M\) with
\Cref{prop:cg-normalized-schur-action} shows that the last line of
\eqref{eq:cg-semialgebraic-schur} is equality of marginal precision and source,
and for nonsingular Gaussian laws this is equality of laws, since covariance
determines precision and \(\mathsf h=\Lambda\mu\) then determines the source,
while the explicit densities give the converse. The admissions have polynomial
encodings. A symmetric matrix is SPD exactly when its leading principal minors
are positive: completing the square in the first coordinate gives a positive
first pivot and a Schur block whose leading minors are the subsequent leading
minors divided by that pivot, and induction proves the claim. A symmetric
matrix is PSD exactly when it equals \(TT^\top\) for some \(T\), by orthogonal
diagonalization and \(v^\top TT^\top v=\|T^\top v\|^2\). Invertibility of a link
is \((\det\Theta)^2>0\). A denominator \(q\) is replaced by a variable \(u\) with
\(uq=1\), which keeps its nonzero domain; multiplying an equation by a nonzero
denominator is valid, whereas an inequality needs this variable or a split into
the two sign cases. An algebraic coefficient with defining polynomial \(f\) and
isolating interval \((a_0,b_0)\) is replaced by an existentially quantified
variable \(v\) with \(f(v)=0\) and \(a_0<v<b_0\), and clearing the rational
denominators leaves integer coefficients; the equivalence holds at every real
parameter because the constant is the unique root in its interval. Finally, the
minimum in \eqref{eq:cg-semialgebraic-schur-positive} is attained at
\(u_{\min}\), and its minimizing vector has nonzero retained part.
\(\square\)

\theoremheading{Exact finite semialgebraic Kron classification}{thm:cg-semialgebraic-kron-classification}
Under \Cref{def:cg-semialgebraic-description}, there is an algorithm producing
a quantifier-free formula \(\mathrm{Good}(x)\) for exactly the admissible fine
parameters whose exact marginal has an admitted retained presentation, and the
failure locus, and deciding universal closure and nonemptiness:
\begin{equation}
\begin{aligned}
\mathrm{Good}(x)&\ \Longleftrightarrow\
 \mathrm{Adm}(x)\wedge\exists y\bigl[\mathrm{Rep}(x,y)\bigr],\\
\mathrm{Fail}(x)&=\mathrm{Adm}(x)\wedge\neg\mathrm{Good}(x),\\
&\forall x\bigl(\mathrm{Adm}(x)\Rightarrow\exists y\bigl[\mathrm{Rep}(x,y)\bigr]\bigr),
\qquad
\exists x\bigl[\mathrm{Adm}(x)\bigr].
\end{aligned}
\label{eq:cg-semialgebraic-classification}
\end{equation}
The witnesses include the auxiliary inverse and Gram variables, and
nonemptiness is decided separately, since vacuous closure is not an inhabited
construction. A nonempty failure locus contains a point with real algebraic
coordinates. The raw scalar is carried outside the formula by
\begin{equation}
\mathfrak c_{\mathrm c}
=\mathfrak c+\frac12\log\det M
-\frac{\dim z_{\mathrm h}}{2}\log(2\pi)
-\frac12\mathsf h_{\mathrm h}^\top M^{-1}\mathsf h_{\mathrm h},
\label{eq:cg-semialgebraic-raw-scalar}
\end{equation}
and the observed experiment of a linear Gaussian sensor with algebraic matrix
data and SPD noise is included by keeping the observation coordinates among the
retained ones. For a fixed finite acyclic diagram of retained coordinate sets,
with one tuple \(y_v\) at each object and all arrow relations, path identities,
and structural conditions collected in one formula
\(\mathrm{Rep}_{\mathrm{diag}}\), the same algorithm decides
\begin{equation}
\forall x\Bigl(\mathrm{Adm}(x)\Rightarrow
\exists (y_v)_{v}
\mathrm{Rep}_{\mathrm{diag}}\bigl(x,(y_v)_{v}\bigr)\Bigr).
\label{eq:cg-semialgebraic-diagram}
\end{equation}
For a fixed finite list of rational operations on one declared domain, totality
and universal one-step domain invariance imply admission of every finite
composition. \status{ESTABLISHED}

\noindent\emph{Proof.}
By \Cref{lem:cg-semialgebraic-encoding}, \(\mathrm{Adm}\) and \(\mathrm{Rep}\) are
finite integer-coefficient formulas, and \(\mathrm{Rep}(x,y)\) holds exactly when
\(y\) is an admitted presentation of the exact marginal, so
\Cref{prop:cg-real-quantifier-elimination} produces \(\mathrm{Good}\) and
\(\mathrm{Fail}\) and decides the two closed sentences. The existential sentence
\(\exists x\bigl[\mathrm{Fail}(x)\bigr]\) has the same truth value over \(\R\) and over the
real algebraic numbers by the same proposition, so a nonempty failure locus has
an algebraic point, which enumeration of finite tuples of algebraic numbers with
exact sign evaluation of \(\mathrm{Fail}\) eventually finds; the point need not
be rational or small. The raw scalar is the third line of
\eqref{eq:cg-kron-complete-schur}; it is a transcendental identity rather than an
ordered-field constraint, and conditions on logarithms, exponentials, or
normalizers enter the decision only after a valid polynomial reduction has been
supplied. For a linear sensor with SPD noise covariance \(\Sigma_o\), the joint
of latent and observation coordinates is a proper Gaussian with rational natural
parameters, its quadratic being the prior quadratic plus
\((o-\mathsf Cz)^\top\Sigma_o^{-1}(o-\mathsf Cz)\), so the classification applies to
it; pushing the whole observed joint preserves its evidence and the
conditional-KL identity of \Cref{sec:cg-kl-recovery}. For a diagram, one formula
with shared variables is decided, whereas separate witnesses for each arrow
would not give a compatible value at a shared object; coordinate
marginalization composes by the equality of iterated Gaussian integrals, sources
and raw scalars included. The composition statement is induction on the number
of operations, undefined denominators counting as failures of totality.
\(\square\)

\paragraph{Bounding an unmarked positive presentation.}
An unrestricted number of independently scalable unmarked summands can also be
reduced to finite input. Let \(\R^n=\bigoplus_{i\in V}\R^{d_i}\) with coordinate
projections \(\operatorname{pr}_i\). An anchor contributes
\(\operatorname{pr}_i^\top A_i\operatorname{pr}_i\) with \(A_i\succeq0\), and for
nodes of equal dimension a link contributes
\begin{equation}
\mathsf Q_e=(\operatorname{pr}_i-\Theta_e\operatorname{pr}_j)^\top W_e(\operatorname{pr}_i-\Theta_e\operatorname{pr}_j),
\qquad W_e\succ0,\quad\det\Theta_e\neq0.
\label{eq:cg-finite-pair-generator}
\end{equation}
Allowed endpoint types form a fixed finite list, every generator may be
independently rescaled by a positive number, and unmarked copies may be deleted.

\propositionheading{Finite positive pair-presentation bound}{prop:cg-finite-pair-presentation-bound}
If \(Q\in\Sym^n\) is a finite sum of the anchors and links just declared, it is
such a sum with at most \(n(n+1)/2\) nonzero summands, and with unrestricted PSD
anchors it has at most one anchor at each node and no self-loops; every weight
stays SPD and every link invertible. Membership in this unmarked presentation
class is equivalent to one finite disjunction of existential polynomial
systems, whose projection onto the entries of \(Q\) is semialgebraic and
decidable at exact algebraic inputs. The scaling and deletion hypotheses are
material: in one scalar dimension, generators restricted to the value \(1\)
represent \(2\) with two copies but not within the bound \(1\). The class need
not be closed, and replacing SPD weights by PSD ones changes membership: for two
two-dimensional nodes and \(u=(1,0)^\top\), the matrix
\begin{equation}
Q_\partial=\begin{pmatrix}u\\-u\end{pmatrix}
\begin{pmatrix}u\\-u\end{pmatrix}^{\top}
\label{eq:cg-finite-pair-strict-boundary}
\end{equation}
is a limit of single links with \(\Theta=I\) and
\(W=\operatorname{diag}(1,\delta)\), \(\delta\downarrow0\), but is not a sum of
admitted generators. \status{ESTABLISHED}

\noindent\emph{Proof.}
A loop is the PSD anchor
\(\operatorname{pr}_i^\top(I-\Theta_e)^\top W_e(I-\Theta_e)\operatorname{pr}_i\).
Absorb unmarked loops into anchors, sum the anchors at each node, and delete zero
terms. Reversing a non-loop orientation replaces \((\Theta,W)\) by
\((\Theta^{-1},\Theta^\top W\Theta)\) and preserves the residual quadratic, so
one orientation suffices when the link domain is the whole general linear group;
a restricted domain keeps its finite orientation types. Write the remaining sum
as \(Q=\sum_{r=1}^mt_r\mathsf Q_r\), initially with \(t_r=1\). If
\(m>n(n+1)/2=\dim\Sym^n\), choose a nonzero dependence
\(\sum_r\xi_r\mathsf Q_r=0\); every nonzero PSD matrix has positive trace, so the
\(\xi_r\) have both signs, and with
\begin{equation}
\bar t=\min_{\xi_r>0}\frac{t_r}{\xi_r},
\qquad t'_r=t_r-\bar t\xi_r,
\label{eq:cg-finite-pair-support-reduction}
\end{equation}
all \(t'_r\geq0\), one vanishes, and \(\sum_rt'_r\mathsf Q_r=Q\). Repetition
terminates, and each surviving coefficient is absorbed into its \(A_i\) or
\(W_e\) without changing the admissions. This is Carath\'eodory's theorem
\citep[Thm.~17.1]{Rockafellar1970} applied to \(Q/\Tr Q\) in the convex hull of
the normalized generators \(\mathsf Q_r/\Tr\mathsf Q_r\), which lie in an affine
hyperplane of dimension \(n(n+1)/2-1\). Only finitely many endpoint and anchor
patterns have at most \(n(n+1)/2\) terms; for each, with symmetric anchor and
weight variables and unrestricted link variables, impose
\begin{equation}
\begin{aligned}
Q&=\sum_i\operatorname{pr}_i^\top A_i\operatorname{pr}_i+\sum_e
(\operatorname{pr}_i-\Theta_e\operatorname{pr}_j)^\top W_e(\operatorname{pr}_i-\Theta_e\operatorname{pr}_j),\\
A_i&=F_iF_i^\top,\qquad
\Tr(F_iF_i^\top)>0,\\
&\text{all leading principal minors of }W_e\text{ positive},
\qquad(\det\Theta_e)^2>0,
\end{aligned}
\label{eq:cg-finite-pair-polynomial-system}
\end{equation}
with the factors \(F_i\) existential. By the encodings of
\Cref{lem:cg-semialgebraic-encoding} every solution is a permitted presentation,
and the reduction puts every finite presentation into one listed pattern, so the
disjunction is necessary and sufficient; quantifier elimination makes its
projection semialgebraic. The scalar example has \(\dim\Sym^1=1\). The matrix
\(Q_\partial\) has rank one and a nonzero cross block, while every non-loop
generator on those nodes has rank two, because \((I,-\Theta)\) has full row rank
and \(W\succ0\); PSD summands of \(Q_\partial\) vanish on its kernel, so their
ranges lie in its one-dimensional range, which excludes every non-loop term, and
anchors and loops have zero cross blocks. The links with
\(W=\operatorname{diag}(1,\delta)\) converge to \(Q_\partial\). \(\square\)

If \(Q\succ0\), repacking \(Q\) while carrying \(\mathsf h\), \(\mathfrak c\), and
the reference unchanged keeps the complete Gaussian action, its mass, and its
normalized law, whereas independently normalized local factors introduce
determinant scalars that must be compensated when raw equality is required.
Individually marked edges, their inventories, and prescribed loop products
cannot be discarded; such problems keep their finite inventory in
\Cref{thm:cg-semialgebraic-kron-classification}, and the proper example of
\Cref{prop:cg-kron-positive-pair-obstruction} excludes even the larger class of
PSD two-node forms.

\propositionheading{Measurable selection and structural obstructions}{prop:cg-kron-selector-boundary}
For the good locus of \Cref{thm:cg-semialgebraic-kron-classification} there is
a semialgebraic, hence Borel, selector \(s\) with \(\mathrm{Rep}(x,s(x))\) for every
\(x\) in the locus. Neither a continuous selector nor a gauge-equivariant
structural selector follows from these hypotheses, even when every accompanying
Gaussian law is proper and identical. For the parameter
\((\cos t,\sin t)\in S^1\) and the polynomial condition \(O\in\mathrm{SO}(2)\),
\(O^2=R_t\), each fiber has two points and no continuous selection exists,
while the state law may stay \(\mathcal N(0,I_2)\). For the input covariance
\(I_2\) and a factor \(O\) with \(OO^\top=I_2\), acted on by
\(O\mapsto gO\), \(g\in\mathrm{SO}(2)\), no equivariant selector exists.
\status{ESTABLISHED}

\noindent\emph{Proof.}
The semialgebraic sets are the definable sets of the real field, an o-minimal
expansion of an ordered group, which has definable choice \citep[Ch.~6, pp.~93--94]{vandenDries1998}: the
semialgebraic family of fibers \(\{y:\mathrm{Rep}(x,y)\}\), \(x\) in the good
locus, has a semialgebraic selector. A semialgebraic map is Borel, since
preimages of rational open boxes are semialgebraic and such boxes form a
countable base. The solutions of \(O^2=R_t\) in \(\mathrm{SO}(2)\) are
\(\pm R_{t/2}\), so along \(t\in[0,2\pi]\) a continuous selection is
\(\sigma R_{t/2}\) with \(\sigma\in\{I,-I\}\) continuous, hence constant, and
then its value at \(2\pi\) is minus its value at \(0\) although the parameter has
returned. The covariance \(I_2\) is fixed by every \(g\in\mathrm{SO}(2)\), so an
equivariant selector would satisfy \(O=gO\) for every \(g\), and \(g=-I_2\)
forces \(O=0\), contradicting \(OO^\top=I_2\). \(\square\)

The selected tuple is a measurable presentation of the exact marginal; it is not
a smooth frame, an equivariant section, or a construction compatible with an
unstated global composition law. Finite diagram compatibility is handled by
\eqref{eq:cg-semialgebraic-diagram} with its shared variables.

\propositionheading{Obstruction for computably enumerated domains}{prop:cg-kron-enumerator-obstruction}
There is no terminating algorithm which, from arbitrary code for a total
computable sequence of rational parameters, decides whether every member of the
sequence has a positive pairwise coarse presentation. This remains true when the
topology and node dimensions are fixed, every fine law is proper, and the range
is promised to lie in one fixed two-point set; the obstruction concerns the
input description, although every such range is finite and semialgebraic as a
set. \status{ESTABLISHED}

\noindent\emph{Proof.}
Use the four-agent, two-dimensional star of
\Cref{prop:cg-kron-positive-pair-obstruction}. Let \(x_{\mathrm{bad}}\) be its
rational edge data with retained anchors \(\delta I\), \(\delta=1/20000\), whose
proper marginal has no positive one-node or two-node presentation. Let
\(x_{\mathrm{good}}\) have the same topology and anchors and all three star
weights equal to \(I\); the hidden pivot is \(3I\), the retained blocks are
\(2I/3\) on the diagonal and \(-I/3\) off it, and the complete retained graph with
weights \(I/3\) together with the anchors is an exact positive presentation. For
a program \(\mathtt p\) define
\begin{equation}
x_k(\mathtt p)=
\begin{cases}
x_{\mathrm{bad}},&\mathtt p\text{ halts within }k\text{ steps},\\
x_{\mathrm{good}},&\text{otherwise}.
\end{cases}
\label{eq:cg-kron-enumerator-reduction}
\end{equation}
Computing a term takes \(k\) simulation steps, so this is a total computable
rational sequence with range in \(\{x_{\mathrm{good}},x_{\mathrm{bad}}\}\), and all
its members admit positive pairwise presentations exactly when \(\mathtt p\) never
halts. An algorithm deciding the property would therefore decide halting. No
machine decides whether an arbitrary machine ever prints a given symbol
\citep[Sec.~8]{Turing1937}, and a program that simulates a machine and halts when
the machine first prints that symbol reduces this problem to halting, so no such
algorithm exists. \(\square\)

There is no conflict with \Cref{thm:cg-semialgebraic-kron-classification}: an
explicit semialgebraic formula for the range is different input from code that
enumerates it. The proposition does not exclude a theorem about a particular
infinite family, a nonalgorithmic classification, or exact full-Gaussian
marginalization.

\section{Finite Gaussian quotients with variable support}
\label{sec:cg-variable-support-quotients}

The full-SPD assumption of \Cref{thm:cg-semialgebraic-kron-classification} can
be replaced by an explicit quotient reference, which gives a finite
classification that includes variable ranks.

\definitionheading{Quotient Gaussian data}{def:cg-quotient-reference}
Quotient Gaussian data on \(\R^n\) consist of a symmetric PSD form \(Q\), an SPD
reference metric \(\mathsf G_{\mathrm{ref}}\), a covector
\(\mathsf h\in\operatorname{range}Q\), an affine origin \(z_0\in\R^n\), and a
real scalar \(\mathfrak c\). The representative slice and the action on it are
\begin{equation}
\mathcal S=(\ker Q)^{\perp_{\mathsf G_{\mathrm{ref}}}},\qquad
\mathsf E(z_0+s)=\mathfrak c+\frac12s^\top Qs-\mathsf h^\top s,\quad s\in\mathcal S,
\label{eq:cg-quotient-reference-action}
\end{equation}
with the volume induced by \(\mathsf G_{\mathrm{ref}}\) on \(\mathcal S\),
translated to \(z_0+\mathcal S\), and unit mass on the point in dimension zero.
Since \(Q\) is symmetric, \(\operatorname{range}Q=(\ker Q)^\perp\), so both
energies descend to \(\R^n/\ker Q\). The metric, origin, and volume specify a
law on a representative slice, not a normalization of a singular ambient
Lebesgue density; a law on the discarded null coordinates would define a
different completion.

\propositionheading{Polynomial quotient covariance}{prop:cg-quotient-covariance-polynomial}
For every such \(Q\) and \(\mathsf G_{\mathrm{ref}}\) there is exactly one
symmetric \(\Sigma_Q\) with
\begin{equation}
Q\Sigma_QQ=Q,\qquad \Sigma_QQ\Sigma_Q=\Sigma_Q,\qquad
\mathsf G_{\mathrm{ref}}\Sigma_QQ=Q\Sigma_Q\mathsf G_{\mathrm{ref}}.
\label{eq:cg-quotient-covariance-graph}
\end{equation}
With the whitened forms
\begin{equation}
\widehat Q=\mathsf G_{\mathrm{ref}}^{-1/2}Q\mathsf G_{\mathrm{ref}}^{-1/2},\qquad
\widehat\Sigma=\mathsf G_{\mathrm{ref}}^{1/2}\Sigma_Q\mathsf G_{\mathrm{ref}}^{1/2},
\label{eq:cg-quotient-whitened-graph}
\end{equation}
it is
\begin{equation}
\Sigma_Q=\mathsf G_{\mathrm{ref}}^{-1/2}\widehat Q^+\mathsf G_{\mathrm{ref}}^{-1/2},
\label{eq:cg-quotient-metric-inverse}
\end{equation}
where \(\widehat Q^+\) inverts the positive eigenvalues and vanishes on the
kernel. \(\Sigma_Q\) is PSD with range \(\mathcal S\), and
\(\Pi_Q=\Sigma_QQ\) is the \(\mathsf G_{\mathrm{ref}}\)-orthogonal projection onto
\(\mathcal S\) with kernel \(\ker Q\). The equations
\eqref{eq:cg-quotient-covariance-graph} are the defining equations of the
weighted Moore--Penrose inverse of \(Q\) with weights
\(\mathsf G_{\mathrm{ref}}^{-1}\) and \(\mathsf G_{\mathrm{ref}}\), because for
symmetric \(\Sigma_Q\) and \(Q\) both weighted symmetry conditions reduce to the
third equation, and they specify the covariance at every rank.
\status{ESTABLISHED}

\noindent\emph{Proof.}
The SPD square root of \(\mathsf G_{\mathrm{ref}}\) exists by orthogonal
diagonalization, and in the variables \eqref{eq:cg-quotient-whitened-graph} the
equations become \(\widehat Q\widehat\Sigma\widehat Q=\widehat Q\),
\(\widehat\Sigma\widehat Q\widehat\Sigma=\widehat\Sigma\), and
\(\widehat\Sigma\widehat Q=\widehat Q\widehat\Sigma\), with \(\widehat\Sigma\)
symmetric. Decompose into the positive eigenspace of \(\widehat Q\) and its
kernel; commutation makes the off-diagonal blocks vanish, the first equation
gives \(\widehat\Sigma=\widehat Q^{-1}\) on the positive block, and the second
gives \(\widehat\Sigma=0\) on the kernel. This proves uniqueness,
\eqref{eq:cg-quotient-metric-inverse}, and positivity. In whitened coordinates
\(\widehat\Sigma\widehat Q\) is the Euclidean orthogonal projection onto
\(\operatorname{range}\widehat Q\), and since
\(\ker\widehat Q=\mathsf G_{\mathrm{ref}}^{1/2}\ker Q\), undoing the isometry
\(s\mapsto\mathsf G_{\mathrm{ref}}^{1/2}s\) gives the stated projection and range.
For the weighted conditions, multiplying
\((\mathsf G_{\mathrm{ref}}^{-1}Q\Sigma_Q)^\top=\mathsf G_{\mathrm{ref}}^{-1}Q\Sigma_Q\)
on both sides by \(\mathsf G_{\mathrm{ref}}\) gives the third equation, and
\((\mathsf G_{\mathrm{ref}}\Sigma_QQ)^\top=\mathsf G_{\mathrm{ref}}\Sigma_QQ\) is the
third equation itself. \(\square\)

The PSD and SPD admissions have the finite encodings of
\Cref{lem:cg-semialgebraic-encoding}, and \(\mathsf h\in\operatorname{range}Q\)
is the existential equation \(Qv=\mathsf h\); the square root and the spectral
inverse fix the meaning of the equations and are not operations of the
polynomial input.

\theoremheading{Normalized quotient Gaussian law}{thm:cg-quotient-gaussian-law}
Let \(r=\operatorname{rank}Q\) and let \(\det_+(\widehat Q)\) be the product of the
positive eigenvalues of \(\widehat Q\), with empty product one. Relative to the
volume on \(z_0+\mathcal S\), the action \eqref{eq:cg-quotient-reference-action}
has positive finite mass and defines the Gaussian probability measure with
ambient mean \(\mu\), covariance \(\Sigma_Q\), and characteristic function
\begin{equation}
\begin{aligned}
Z_Q&=e^{-\mathfrak c}(2\pi)^{r/2}
       \det\nolimits_+(\widehat Q)^{-1/2}
       \exp\Bigl(\frac12\mathsf h^\top\Sigma_Q\mathsf h\Bigr),\\
\mu&=z_0+\Sigma_Q\mathsf h,\\
\widehat{\mathbb P}_Q(\xi)
 &=\exp\Bigl(i\xi^\top\mu-\frac12\xi^\top\Sigma_Q\xi\Bigr).
\end{aligned}
\label{eq:cg-quotient-gaussian-law}
\end{equation}
For \(Q=0\) these give \(\mathsf h=\Sigma_Q=0\), \(Z_Q=e^{-\mathfrak c}\), and the point
mass at \(z_0\). Under an invertible passive change \(z'=gz\) with all data
transformed,
\begin{equation}
Q'=g^{-\top}Qg^{-1},\quad
\mathsf G_{\mathrm{ref}}'=g^{-\top}\mathsf G_{\mathrm{ref}}g^{-1},\quad
\mathsf h'=g^{-\top}\mathsf h,\quad z_0'=gz_0,
\label{eq:cg-quotient-gauge-data}
\end{equation}
one has \(\Sigma_Q'=g\Sigma_Qg^\top\), \(\mu'=g\mu\), and
\(\mathbb P'=g_\#\mathbb P\), with \(Z_Q\) and \(\mathfrak c\) unchanged. The
reference metric must transform with the data: for \(Q=\operatorname{diag}(1,0)\),
\(\mathsf G_{\mathrm{ref}}=I_2\), and
\(g=\begin{psmallmatrix}1&1\\0&1\end{psmallmatrix}\), the pushed covariance is
\(\operatorname{diag}(1,0)\), while resetting the transformed metric to \(I_2\)
gives
\begin{equation}
Q'=\begin{pmatrix}1&-1\\-1&1\end{pmatrix},
\qquad
(Q')^+_{\mathrm{Euclidean}}
=\frac14\begin{pmatrix}1&-1\\-1&1\end{pmatrix},
\label{eq:cg-quotient-reference-counterexample}
\end{equation}
a different slice and law. \status{ESTABLISHED}

\noindent\emph{Proof.}
In the coordinates \(y=\mathsf G_{\mathrm{ref}}^{1/2}s\) the slice becomes
\(\operatorname{range}\widehat Q\) and its volume becomes Euclidean volume; the
restriction of \(\widehat Q\) there is positive definite, and the source
\(\mathsf G_{\mathrm{ref}}^{-1/2}\mathsf h\) lies in that range because it
annihilates \(\ker\widehat Q\). Diagonalizing and completing \(r\) scalar squares
gives the first line of \eqref{eq:cg-quotient-gaussian-law}; transforming the
mean and covariance back gives \(\Sigma_Q\mathsf h\) and \(\Sigma_Q\), and
translation adds \(z_0\). The same one-dimensional integrations give the
characteristic function, and in dimension zero the integral is over one point,
so \(0<Z_Q<\infty\) in all cases. Substituting \eqref{eq:cg-quotient-gauge-data}
into \eqref{eq:cg-quotient-covariance-graph} and using uniqueness gives
\(\Sigma_Q'=g\Sigma_Qg^\top\); \(g\) is an isometry between the two representative
metrics, so it pushes their volume measures and leaves \(Z_Q\) unchanged. The
counterexample is direct computation. \(\square\)

\propositionheading{Pushforwards of quotient Gaussian laws}{prop:cg-quotient-pushforward}
(i) Every affine map \(z\mapsto\mathsf F z+\mathfrak b\), including rectangular or
rank-deficient \(\mathsf F\), pushes the law of
\Cref{thm:cg-quotient-gaussian-law} to the Gaussian with
\begin{equation}
\mu_{\mathrm{push}}=\mathsf F(z_0+\Sigma_Q\mathsf h)+\mathfrak b,\qquad
\Sigma_{\mathrm{push}}=\mathsf F\Sigma_Q\mathsf F^\top .
\label{eq:cg-quotient-affine-push}
\end{equation}
Two Gaussian measures on one finite ambient space, singular or not, agree
exactly when their means and covariances agree, so target data
\((Q_{\mathrm c},\mathsf G_{\mathrm{ref,c}},\mathsf h_{\mathrm c},z_{0,\mathrm c})\) with
covariance \(\Sigma_{\mathrm c}\) represent this push exactly when
\begin{equation}
\Sigma_{\mathrm c}=\mathsf F\Sigma_Q\mathsf F^\top,\qquad
z_{0,\mathrm c}+\Sigma_{\mathrm c}\mathsf h_{\mathrm c}=\mathsf F(z_0+\Sigma_Q\mathsf h)+\mathfrak b,
\label{eq:cg-quotient-ambient-law-equality}
\end{equation}
as equalities of probability measures, supports included.
(ii) A map \([z]\mapsto[\mathsf F z]\) between the energy quotients is well defined
exactly when \(\mathsf F\ker Q\subseteq\ker Q_{\mathrm c}\), that is,
\begin{equation}
Q_{\mathrm c}\mathsf F(I-\Sigma_QQ)=0.
\label{eq:cg-quotient-map-descent}
\end{equation}
With \(\Pi_{\mathrm c}=\Sigma_{\mathrm c}Q_{\mathrm c}\) and a declared quotient
offset \(\mathfrak b_0\), the target representative map
\(z_0+s\mapsto z_{0,\mathrm c}+\Pi_{\mathrm c}(\mathsf F s+\mathfrak b_0)\) pushes the law
to the target law exactly when
\begin{equation}
\Sigma_{\mathrm c}=\Pi_{\mathrm c}\mathsf F\Sigma_Q\mathsf F^\top\Pi_{\mathrm c}^\top,\qquad
\Sigma_{\mathrm c}\mathsf h_{\mathrm c}=\Pi_{\mathrm c}(\mathsf F\Sigma_Q\mathsf h+\mathfrak b_0);
\label{eq:cg-quotient-intrinsic-law-equality}
\end{equation}
choosing a map on one slice alone does not prove descent.
(iii) For three scalar nodes with two unit edges, no anchors, zero source, and
\(\mathsf G_{\mathrm{ref}}=I_3\),
\begin{equation}
Q_{\mathrm f}=\begin{pmatrix}
1&-1&0\\-1&2&-1\\0&-1&1
\end{pmatrix},\qquad
\Sigma_{\mathrm f}=\frac19\begin{pmatrix}
5&-1&-4\\-1&2&-1\\-4&-1&5
\end{pmatrix},
\label{eq:cg-quotient-chain-data}
\end{equation}
and \(\mathsf F\) selecting the first and third coordinates, with
\(\mathsf G_{\mathrm{ref,c}}=I_2\),
\begin{equation}
Q_{\mathrm c}=\Sigma_{\mathrm c}=\Pi_{\mathrm c}
=\frac12\begin{pmatrix}1&-1\\-1&1\end{pmatrix},\qquad
\mathsf F\Sigma_{\mathrm f}\mathsf F^\top
=\frac19\begin{pmatrix}5&-4\\-4&5\end{pmatrix};
\label{eq:cg-quotient-chain-push}
\end{equation}
the map descends and
\(\Pi_{\mathrm c}\mathsf F\Sigma_{\mathrm f}\mathsf F^\top\Pi_{\mathrm c}^\top=\Sigma_{\mathrm c}\),
so the coarse edge of weight \(1/2\) is an exact quotient push with zero sources
and offsets, whereas the unprojected endpoint covariance differs from
\(\Sigma_{\mathrm c}\). \status{ESTABLISHED}

\noindent\emph{Proof.}
(i) Substitute \(\mathsf F^\top\xi\) in the characteristic function and multiply by
\(e^{i\xi^\top\mathfrak b}\). Equal means and covariances give equal
characteristic functions, and the converse follows from finite first and second
moments. (ii) Replacing \(z\) by \(z+k\) with \(k\in\ker Q\) shows necessity of
\(\mathsf F\ker Q\subseteq\ker Q_{\mathrm c}\), and the same calculation shows
sufficiency; \(I-\Sigma_QQ\) projects onto \(\ker Q\) by
\Cref{prop:cg-quotient-covariance-polynomial}, which gives
\eqref{eq:cg-quotient-map-descent}. The representative map is affine, so (i)
applied to it gives \eqref{eq:cg-quotient-intrinsic-law-equality}, the target
mean being \(z_{0,\mathrm c}+\Sigma_{\mathrm c}\mathsf h_{\mathrm c}\). (iii)
Multiplication gives \(Q_{\mathrm f}\Sigma_{\mathrm f}=I_3-\mathbf1\mathbf1^\top/3\)
and the covariance equations; the fine common-translation kernel maps into the
coarse one, and direct multiplication gives the projected covariance.
Equivalently, the fine residuals \(z_1-z_2\) and \(z_2-z_3\) are independent
standard normals, and their sum has variance two. \(\square\)

A target metric or projection that varies with a model parameter does not by
itself define one parameter-independent statistical channel, so Fisher
contraction for a family requires that admission separately.

\theoremheading{Finite variable-support Gaussian classification}{thm:cg-variable-support-classification}
Fix finite ambient dimensions, finitely many declared affine or intrinsic
quotient maps, and finite semialgebraic input and target predicates over
effectively represented real algebraic coefficients, with rational matrix
expressions on explicit nonzero-denominator domains. At each object require
\(Q\succeq0\), \(\mathsf G_{\mathrm{ref}}\succ0\), and
\(\mathsf h\in\operatorname{range}Q\), with reference and origin as in
\Cref{def:cg-quotient-reference}; ranks and supports may vary with the
parameters. Then an algorithm constructs the exact good and failure loci for
quotient-Gaussian representability,
\begin{equation}
\begin{aligned}
\mathrm{Good}_{\mathrm{quot}}(x)&\Longleftrightarrow
\mathrm{Adm}(x)\wedge\exists y\bigl[\mathrm{Rep}_{\mathrm{quot}}(x,y)\bigr],\\
\mathrm{Fail}_{\mathrm{quot}}(x)&=\mathrm{Adm}(x)\wedge\neg\mathrm{Good}_{\mathrm{quot}}(x),
\end{aligned}
\label{eq:cg-variable-support-good-locus}
\end{equation}
and decides universal closure, separate nonemptiness, and any finite
simultaneous diagram with its supplied polynomial structural compatibilities;
for affine maps it uses \eqref{eq:cg-quotient-ambient-law-equality}, and for
intrinsic maps both \eqref{eq:cg-quotient-map-descent} and
\eqref{eq:cg-quotient-intrinsic-law-equality}. The asserted equivalence is
equality of the normalized measures on their representative spaces. The raw
finite measures are carried by the analytic normalizer at each object: with
\begin{equation}
\begin{aligned}
\mathsf A_i^{\mathrm{mass}}
&=\frac{r_i}{2}\log(2\pi)
-\frac12\log\det\nolimits_+(\widehat Q_i)
+\frac12\mathsf h_i^\top\Sigma_{Q_i}\mathsf h_i,\\
\mathfrak c_{\mathrm c}
&=\mathfrak c_{\mathrm f}
+\mathsf A_{\mathrm c}^{\mathrm{mass}}
-\mathsf A_{\mathrm f}^{\mathrm{mass}},
\end{aligned}
\label{eq:cg-quotient-raw-scalar}
\end{equation}
a free target scalar gives \(Z_{\mathrm c}=Z_{\mathrm f}\), and equality of the
normalized laws then gives equality of the pushed raw finite measures. An
observed experiment is included by keeping its observation coordinates in the
finite joint and requiring the representative map to leave them unchanged,
which preserves its observation marginal, singular or not.
\status{ESTABLISHED}

\noindent\emph{Proof.}
Introduce a symmetric covariance variable at each object and impose
\eqref{eq:cg-quotient-covariance-graph}, \(Qv=\mathsf h\), and the PSD and SPD
admissions; uniqueness in \Cref{prop:cg-quotient-covariance-polynomial} ensures
that an existential covariance cannot select a different law. Denominators are
handled as in \Cref{lem:cg-semialgebraic-encoding}, so the map equations are
finite polynomial relations, and by \Cref{prop:cg-quotient-pushforward} they are
necessary and sufficient for the normalized pushforward, descent included when
an intrinsic map is requested. A prescribed rank \(r>0\) is the vanishing of
all \((r+1)\)-minors and positivity of the sum of squares of the \(r\)-minors,
rank zero is the vanishing of all entries, and the finitely many ranks form a
finite disjunction; \(\operatorname{range}U\subseteq\operatorname{range}U'\) is
the existential equation \(U=U'C\), and affine support relations add membership
of the displacement. \Cref{prop:cg-real-quantifier-elimination} then produces
\eqref{eq:cg-variable-support-good-locus} and decides the closed sentences. For
a diagram, one target tuple and its unique covariance at each object carry all
arrow equations simultaneously, with descent and the declared identities of the
induced quotient maps for quotient arrows; equality modulo a target null
direction is encoded as such, and composable representative maps compose their
moment transformations by affine substitution. No nonsingular inverse or common
support is assumed. The raw-scalar choice gives \(Z_{\mathrm c}=Z_{\mathrm f}\) by
\eqref{eq:cg-quotient-gaussian-law}, and a common positive mass factor cancels
from a normalized law, so raw equality is a separate requirement that this
choice meets. \(\square\)

A preferred positive pair factorization is an extra target condition,
expressible by \eqref{eq:cg-finite-pair-polynomial-system} when its hypotheses
hold, or through a supplied marked inventory; variable support does not relax
positive definiteness of edge weights, erase holonomy marks, or enlarge a
prescribed target domain. Logarithms, positive pseudodeterminants, and
rank-dependent normalizers enter the polynomial decision only through such a
reduction.

\propositionheading{Rank change obstructs statistical regularity}{prop:cg-rank-change-statistical-obstruction}
Finite semialgebraic feasibility across ranks does not imply total-variation
continuity, differentiability in quadratic mean, or a smooth positive Fisher
geometry, even for a smooth PSD coefficient path whose quotient law is
normalized at every parameter. Even the smooth covariance path
\(\operatorname{diag}(1,t^2)\) has total-variation distance one from its
line-supported law at \(t=0\), although these laws converge weakly.
\status{ESTABLISHED}

\noindent\emph{Proof.}
Take
\begin{equation}
Q(t)=\operatorname{diag}(1,t^2),\qquad
\mathsf G_{\mathrm{ref}}=I_2,\qquad \mathsf h=z_0=0.
\label{eq:cg-rank-change-example}
\end{equation}
For \(t\neq0\) the covariance is \(\operatorname{diag}(1,t^{-2})\) and the
support is \(\R^2\); at zero the covariance is \(\operatorname{diag}(1,0)\) and the
support is the first coordinate axis, which has probability zero for every
\(t\neq0\) and one at zero, so with
\(\|\mathbb P-\mathbb P'\|_{\mathrm{TV}}=\sup_B|\mathbb P(B)-\mathbb P'(B)|\) the
distance is exactly one. For every finite \(R_0\),
\begin{equation}
\mathbb P_t(|X_2|\leq R_0)
=\mathbb P(|\varsigma|\leq R_0|t|)\longrightarrow0,
\qquad \varsigma\sim\mathcal N(0,1),
\label{eq:cg-rank-change-escape}
\end{equation}
so there is no weak probability limit as \(t\to0\): the bounded continuous
positive function \(e^{-X_2^2}\) has expectation \(|t|/\sqrt{t^2+2}\to0\),
whereas its integral against any probability measure on \(\R^2\) is positive.
On \(t>0\) the second-coordinate log density has derivative \(1/t-tX_2^2\), with
mean zero and variance \(2/t^2\), the Fisher information. Two-dimensional
Lebesgue measure plus the one-dimensional measure on the axis dominates all
these laws. Differentiability in quadratic mean at zero would require a
square-integrable score \(\mathbf s_0\) with
\(\|\sqrt{p_t}-\sqrt{p_0}-\frac t2\mathbf s_0\sqrt{p_0}\|_2=o(|t|)\); the triangle
inequality would give \(\|\sqrt{p_t}-\sqrt{p_0}\|_2\to0\), and Cauchy--Schwarz
applied to \((\sqrt{p_t}-\sqrt{p_0})(\sqrt{p_t}+\sqrt{p_0})\) would give
\(\|\mathbb P_t-\mathbb P_0\|_{\mathrm{TV}}\leq\|\sqrt{p_t}-\sqrt{p_0}\|_2\to0\),
contradicting the distance-one calculation. For the covariance path
\(\operatorname{diag}(1,t^2)\) the same support argument gives distance one,
and the coupling \((\varsigma_1,t\varsigma_2)\) with dominated convergence gives weak
convergence; its quotient precisions satisfy
\eqref{eq:cg-quotient-covariance-graph}, so weak convergence alone does not
supply a regular statistical geometry. \(\square\)

\openproblemheading{Kron classification beyond finite encodings}{open:cg-kron-classification-beyond-finite}
Finite algebraically encoded full-SPD problems are decided by
\Cref{thm:cg-semialgebraic-kron-classification}, and declared metric quotient
laws with variable rank and support by
\Cref{thm:cg-variable-support-classification}; correlated, nonconvex,
variable-anchor, and branching nonflat domains are included when their
constraints fit the finite input, and the scalar formulas are exact analytic
outputs outside the polynomial formulas. Universal continuity and equivariance
of selectors (\Cref{prop:cg-kron-selector-boundary}), smooth passage between
support strata (\Cref{prop:cg-rank-change-statistical-obstruction}), and a
universal decision procedure for enumerated inputs
(\Cref{prop:cg-kron-enumerator-obstruction}) are refuted. Unproved is any
classification, closed form or algorithmic, of families outside finite
encodings, such as nonalgebraic domains, changing-dimension limits, and
alternative completions of the null coordinates, and any construction with a
global regularity target, such as a smooth, equivariant, or functorial
structural selection. For a stated family, a declaration of its domain and
compatibility conditions together with a proof of the requested invariance,
limit, or global construction would settle it, as would a reduction of the
family to a finite encoding covered by the two theorems.
